%\section{Guessing Profitable elements}
%\label{sec:guess}

%In this section we describe the assignment of elements with relatively large profits 



\section{introduction}
\label{sec:introduction}

The problem of {\em budgeted matroid independent set (BMI)} consists of finding an independent set in a matroid with bounded {\em size} of maximal {\em profit}. TBA...

\section{Preliminaries}
\label{sec:preliminaries}

For the simplicity of the notation, for any set $A$ and some element $a$, let $A+a$, $A-a$ be $A \cup \{a\}$ and $A \setminus \{a\}$, respectively. In addition, for all $k \in \mathbb{N}$ let $[k]$ be the set $\{1,2,\ldots,k\}$. Finally, for any function $f:A \rightarrow B$ and $C \subseteq A$ let $f(C) = \sum_{\ell \in C} f(\ell)$. 

\subsection{Matroids}

Let $E$ be a finite {\em ground set} and let $\cI \subseteq 2^E$ be a non-empty set containing subsets of $E$ called the {\em independent sets} of $E$. Then, $\cm = (E, \cI)$ is a {\em matroid} if the following hold: 

\begin{enumerate}
	\item (The Hereditary Property) For all $A \in \cI$ and $B \subseteq A$ it holds that $B \in \cI$.
	
	\item(The exchange Property) For every $A,B \in \cI$ where $|A| > |B|$, there is $x \in A \setminus B$ such that $B +x \in \cI$. 
\end{enumerate}

Throughout this paper, determining whether a subset of elements forms an independent set of a matroid is assumed to be relied on an {\em oracle} and is considered to take polynomial time in the size of a matroid. Specifically, we assume that there is a polynomial time algorithm that given a matroid $\cm = (E, \cI)$ and a subset of elements $S \subseteq E$, returns "yes" if $S \in \cI$ and "no" otherwise. The next lemma gathers known results which follow from the definition of matroid (see, e.g., \cite{schrijver2003combinatorial} for more details). 


\begin{lemma}
	\label{lem:prel}
	
	The following hold for any matroid $\cm = (E, \cI)$. 
	%
	\begin{enumerate}
		
		\item For all $A,B \in \cI$ where $|A| \geq |B|$, there is $C \subseteq A \setminus B$, $|C| = |A|-|B|$ such that $B \cup C \in \cI$. 
		
		
		\item For every $F \subseteq E$ then $\cm - F = (E \setminus F, \cI')$ is a matroid where $\cI' = \{A \in \cI~|~ A \cap F = \emptyset\}$. 
		
		\item For every $F \in \cI$ then $\cm / F = (E \setminus F, \cI')$ is a matroid where $\cI' = \{A \subseteq E'~|~ A \cup F \in \cI\}$. 
		
		\item For every $q \in \mathbb{N}$ then $\cm^q = (E, \cI^q)$ is a matroid where $\cI^q = \{A \in \cI~|~ |A| \leq q\}$. 
	\end{enumerate}
	\comment{
	\begin{enumerate}
		
		\item For all $A,B \in \cI$ where $|A| \geq |B|$, there is $C \subseteq A \setminus B$, $|C| = |A|-|B|$ such that $B \cup C \in \cI$. 
		
		
		\item For every $F \subseteq E$ then $M - F = (E', \cI')$ is a matroid called the deletion of $F$ in $M$ where $E' = E \setminus F$ and $\cI' = \{A \in \cI~|~ A \cap F = \emptyset\}$. 
		
		\item For every $F \in \cI$ then $M / F = (E', \cI')$ is a matroid called the contraction of $F$ in $M$ where $E' = E \setminus F$ and $\cI' = \{A \subseteq E'~|~ A \cup F \in \cI\}$. 
		
			\item For every $q \in \mathbb{N}$ then $M^q = (E, \cI^q)$ is a matroid called the truncation of $q$ in $M$ where $\cI^q = \{A \in \cI~|~ |A| \leq q\}$. 
	\end{enumerate}
}
	\end{lemma}


\subsection{Matroid polytope}

%with $n$ elements numbered as $\{1,\ldots,n\}$
%A {\em basis} of $M$ is an independent set $B \in \cI$ such that for all $a \in E \setminus B$ it holds that $B+a \notin \cI$. Let $\textsf{bases}(M)$ be the set of all bases of $M$. Given a basis $B$ of $M$, 

Let $\cm = (E, \cI)$ be a matroid. Given $B \in \cI$, the {\em indicator vector} of $B$ is the vector $\bar{I}^B \in \{0,1\}^E$ where for all $i \in B$ and $j \in [n] \setminus B$ we have $\bar{I}^B_i = 1$ and  $\bar{I}^B_j = 0$, respectively. Then, the {\em matroid polytope} of $M$ is the convex hull of the set of indicator vectors of all independent sets of $\cm$: $P_{\cm} = \textsf{conv} \{\bar{I}^B~|~B \in \cI\}$. In addition, a {\em face} of a polytope $P$ is the intersection of P with any closed half-space whose boundary is disjoint from the interior of $P$. The next result is due to \cite{grandoni2010approximation}.

 \begin{lemma}
	\label{lem:integral}
	Let $\cm$ be a matroid and let $F$ be a face of dimension
	$d$ of the matroid polytope $P_{\cm}$. Then any $\bar{x} \in F$ has at most $2d$ non-integral
	components.
\end{lemma}


\subsection{Instance and Solution}

%For all $Q \subseteq E$, let $s(Q) = \sum_{\ell \in Q} s(\ell)$ and $p(Q) = \sum_{\ell \in Q} p(\ell)$. 

Formally, a BMI instance is a tuple $\ck =  (E,\cI,s,p)$, where $(E,\cI)$ is a matroid, $s:E \rightarrow [0,1]$  gives the element sizes and $p:E \rightarrow \mathbb{R}^+$ gives the elements profits. A {\em solution} of $\ck$ is an independent set $X \in \cI$ such that $s(X) \leq 1$; in addition, the {\em profit} of the solution is $p(X)$ and let $\OPT(\ck)$ denote the profit of a solution of $\ck$ with maximal profit. %In this section we prove the following result. 



\section{The Algorithm}
\label{sec:alg}


In this section an EPTAS for BMI is presented. The algorithm is relied on handling elements with high profit separately, where we find a small subset of elements from which the most profitable elements of the solution are chosen. Additional elements with lower profits are then added using a {\em linear program (LP)}. For the remainder of this section, fix a BMI instance $\ck =  (E,\cI,s,p)$ and a parameter $0<\eps<1$, where without the loss of generality, we assume that $\eps^{-1} \in \mathbb{N}$ and that for all $\ell \in E$ it holds that $\{\ell\} \in \cI$ (or this element can be omitted). For the flow of the reading, some of the more technical proofs are deferred to Section~\ref{sec:proofs}. 


In the following, we use rounding on profit of elements with relatively large profits. We start with computing a constant factor approximation of the optimum $\OPT(\ck)$. By the results of Grandoni and Zenklusen, the {\em $k$-budgeted matroid independent set} admits a PTAS \cite{grandoni2010approximation}. Since BMI is a special case of the latter problem with a single budget, using parameter $\eps' = \frac{1}{2}$, we can compute in polynomial time a solution $B^*$ of $\ck$ such that $p(B^*) \geq \frac{\OPT(\ck)}{2}$. Moreover, let $t = \eps^{-1}+1$ and let $h = \eps^t \cdot p(B^*)$. 

We say that an element $\ell \in E$ is {\em profitable} if $p(\ell) \geq h$ and {\em unprofitable} otherwise. Let $H$ denote the set of profitable elements; that is, $H = \{\ell \in E~|~ p(\ell) \geq h\}$. In our algorithm, we round the profits of the profitable elements down to the next multiple of $\eps \cdot h$; note that each elementmay lose only a small fraction of the profit by the above rounding (see Lemma~\ref{lem:J}). Formally, define the {\em profit classes} as $$T = \big\{r \cdot \eps \cdot h ~|~r \in \mathbb{N}_{>0}, r \cdot \eps \cdot h \leq 2 \cdot p(B^*)\big\}$$ which is the set containing the possible profits of profitable elements after the rounding. Now, define $\bar{p}: E \rightarrow \mathbb{R}^+$ as the {\em rounded profit function} such that for all unprofitable elements $\ell \in E \setminus H$ it holds that $\bar{p}(\ell) = p(\ell)$ and for all profitable elements $\ell \in H$ it holds that \begin{equation}
	\label{eq:max}
	\bar{p}(\ell) = \max_{\{y \in T~|~ y ~\leq ~ p(\ell)\}} y.
\end{equation} %Now, let $\cJ = (I,\mathcal{G},s,\bar{p},m,c)$ be the instance where the profits are rounded down as defined above. 


\begin{lemma}
	\label{lem:J}
	For any $A \subseteq E$ it holds that $\bar{p}(A) \geq (1-\eps) \cdot p(A)$. 
\end{lemma} 




Hence, by Lemma~\ref{lem:J} finding a solution with high profit w.r.t. $\bar{p}$ yields a solution with almost the same profit w.r.t. $p$; on the other hand, the number of distinct profits of profitable elements for $\bar{p}$ is bounded by a function of $\eps^{-1}$. Let $q = 2 \cdot \eps^{-t-1}$. By the next result, $q$ is an upper bound on the number of profitable elements in a solution of $\cm$ and on the number of profit classes. 

\begin{lemma}
	\label{q}
	For any solution $A$ of $\cm$ it holds that (i) $|H \cap A| \leq q$ and (ii) $|T| \leq q$. 
\end{lemma}



Our scheme finds a subset of the profitable elements $S$ of size bounded by $q^2$; these are the {\em chosen elements}. By enumerating over all subsets of the chosen elements,  we find $F \subseteq S$ which forms the set of profitable elements returned by the scheme. Then, additional unprofitable elements are added to $F$ using an LP. Surprisingly, despite the small number of chosen elements which depends solely on $\eps$, this simple approach loses only a small portion of the overall profit, in comparison to an optimal solution. 




The construction of $S$ is explained below.  For every profit class $y \in T$, let $E_y = \{\ell \in E~|~p(\ell) = y)\}$ be the subset of elements with profit $y$. For all $y \in T$, we construct a subset of elements $B_y \subseteq E_y$ independently of the remaining profit classes. Recall that the {\em truncation} of $q$ in $M$ is the matroid $M^q = (E,\cI^q)$ where $\cI^q = \{A \in \cI~|~ |A| \leq q\}$ contains all independent sets in $\cm$ with cardinality at most $q$. 

We construct $B_y$ by taking the independent set in $E_y \cap \cI^q$ with maximal number of elements (i.e., a {\em basis}) that has a minimal total size of elements. This can be computed easily by the following greedy approach. Let $E_y  = \{y(1), \ldots, y(n_y)\}$ be non decreasing order of $E_y$ by elementsizes; that is, for all $i,j \in [n_y], i \leq j$ we have $s\left(y(i)\right) \leq s\left(y(j)\right)$. Then, initialize $B_y = \emptyset$ and the next element$y(i) \in E_y$ according to the sorted order is added to $B_y$ if $B_y+y(i)$ forms an independent set in the truncation of $q$ in $M$. Finally, the chosen elements are $S = \bigcup_{y \in T} B_y$ that is the union of the elements chosen for $B_y$ for all profit classes. The algorithm for constructing the chosen elements is reffered to as Algorithm \textsf{Chosen} and the pseudocode of the algorithm is given in Algorithm~\ref{alg:chosen}. \ilan{I dont have time in this iteration by I have a suggestion for improving Algorithm Chosen: to define a minimal basis with at most $q$ elementfor each $E_y$, and in the algorithm just write that we take the union of the minimal bases of all $y \in T$.}


 \begin{algorithm}[h]
	\caption{$\textsf{Chosen}((E,\cI,s,p),q,T)$}
	\label{alg:chosen}
	%Compute the instance $\mathcal{J} = (I,\cG,s,\bar{p},m,c)$. 
	
	\For{$y \in T$}{
		
		Initialize $B_y \leftarrow \emptyset$.
		
		Sort $E_y$ in non decreasing order by element ysizes $y(1), \ldots, y(n_y)$.\label{order}
		
		\For{$i = 1,\ldots,n_y$}{
			
			\If{$B_y + y(i) \in \cI^q$ \label{step:if}}{
				
				Update $B_y \leftarrow B_y+y(i)$
				
			} 
			
		}
		
	}
	
	Return $S = \bigcup_{y \in T} B_y$. 
\end{algorithm}




 
 
\begin{theorem}
	\label{lem:main}
	For all $G \in \cI^q$, where $G \subseteq H$, there is $Z \in \cI^q$,  where $Z \subseteq H$, such that the following hold. 
	
	\begin{enumerate}
		\item$s(Z) \leq s(G)$.
		
		\item For all $y \in T$ it holds that $|E_y \cap Z| = |E_y \cap G|$.
		
		\item$Z \subseteq S$. 
	\end{enumerate} 
\end{theorem} The proof of Theorem~\ref{lem:main} is given in Section~\ref{sec:lemMainProof}. Theorem~\ref{lem:main} is the key component in the analysis of the scheme: the profitable elements of an optimal solution may be replaced by a subset of the chosen elements with similar profit without increasing the total size of elements. However, for the unprofitable elements other techniques are required. 

For adding unprofitable elements to the solution, we use an LP based on a matroid polytope.  Let $F \subseteq H$ be some subset of the profitable elements. In the following, we define an LP whose objective is to find a subset of the unprofitable elements of maximal profit that together with $F$ forms a solution of $\cm$. Since $F$ may prevent an addition of an optimal subset of the unprofitable elements due to the matroid constraint, we restrict the potential unprofitable elements to be of significantly smaller profit in comparison with the elements in $F$. In the analysis, this turns out to be beneficial and decreases the potential profit only slightly by choosing an appropriate $F$. Formally, let \begin{equation}
	\label{E_F}
	E_F = \big\{\ell \in E \setminus H~|~ p(\ell) \leq \eps \cdot \min_{j \in F} p(j)\big\}
\end{equation} be the {\em potential elements} for $F$. The elements in $E_F$ are of profits bounded by $\eps$ times the minimal profit of an elementin $F$. Now, define the $F$-{\em matroid} as $M_F = (E_F,\cI_F)$ where  $\cI_F = \{A \subseteq E_F~|~ A \cup F \in \cI\}$. In other words, the $F$-matroid is the deletion of $E\setminus E_F$ in $M$ combined with the contraction of $F$ in $M$; hence, the $F$-matroid is indeed a matroid by Lemma~\ref{lem:prel} for $F \in \cI$. Finally, define the $F$-LP as the following relaxation of the problem:

 \begin{equation}
 	\label{LP}
 	\begin{aligned}
 		~~~~~ \max\quad        & ~~~\sum_{\ell \in E_F} \bar{x}_{\ell}   \cdot p(\ell)             ~~~~~~~~~~~~~~~~~~~~~                                                \\
 		\textsf{s.t.\quad} & ~~~\sum_{\ell \in E_F} \bar{x}_{\ell}  \cdot s(\ell)         \leq  1-s(F) ~~~~~~~~~~~~~~~~~~~~~           \\  
 		& ~~~\bar{x} \in P_{M_F} \\
 		%& ~~~\bar{x} \in [0,1]^{E_F} 
 	\end{aligned}
 \end{equation}  The $F$-LP maximizes the total profit of a point in the $F$-matroid, with total size of elements at most $1-s(F)$; that is, the residual capacity of a solution without the elements in $F$. In particular, observe that there is no solution to the $F$-LP if $s(F) > 1$ and there is only an all-zero solution in case that $F \notin \cI$.  In addition, observe that a vertex $\bar{\lambda}$ of the $F$-LP is in  a face of dimension $1$ (i.e., an edge) of the $F$-matroid polytope $P_{M_F}$, since the $F$-LP is obtained from $P_{M_F}$ by adding a single linear constraint for bounding the capacity. Hence, by Lemma~\ref{lem:integral} it follows that $\bar{\lambda}$ has at most $2$ non-integral entries. Using the above, we have the required components for designing the following EPTAS for BMI. 




 \begin{algorithm}[h]
	\caption{$\textsf{EPTAS}(\cm = (E,\cI,s,p),\eps)$}
	\label{alg:EPTAS}
	%Compute the instance $\mathcal{J} = (I,\cG,s,\bar{p},m,c)$. 
	
	Construct $q,T$ using a $2$-approximation for BMI. 
	
	Construct the chosen elements $S = \textsf{Chosen}(\cm,q,T)$. 
	
	
	Initialize $A \leftarrow \emptyset$.\label{step:init}
	
	\For{$F \subseteq S$ \label{step:for}}{
		
		
		Find a vertex optimal solution $\bar{\lambda}^F$ of the $F$-LP and construct $C(F)$.\label{step:vertex}
		
			
			\If{$C(F)$ is a solution of $\cm$ and $p\left(C(F)\right) > p(A)$}{
				
				Update $A \leftarrow C(F)$\label{step:update}
				
			} 
			
		
	}
	
	Return $A$. 
\end{algorithm}

Let $S$ be the chosen elements returned by Algorithm \textsf{Chosen} on input $\cm,q,T$. For all subsets $F \subseteq S$, we find a vertex $\bar{\lambda}^F$ of the $F$-LP and round down the entries of $\bar{x}^F$: let $\bar{\gamma}^F \in \{0,1\}^{E_F}$ be the vector where $\bar{\gamma}^F_{\ell} = \floor{\bar{\lambda}^F_{\ell}}$.  In addition, let $\textsf{supp}\left(\bar{\gamma}^F\right) = \{\ell \in E_F~|~ \bar{\gamma}^F_{\ell} = 1\}$ be the {\em support} of $\bar{\gamma}^F$ and let the  {\em solution} of $F$ be the set $C(F) = F \cup \textsf{supp}\left(\bar{\gamma}^F\right)$. The solution of $F$ is called the {\em solution } of $F$ if it forms a solution of $\cm$. Finally, our algorithm chooses among the solutions of all $F \subseteq S$ a solution $C(F^*)$ with maximal total profit. The pseudocode of our scheme is given in Algorithm~\ref{alg:EPTAS}. For the following, let $|\cm|$ be the encoding size of $\cm$. 


\begin{lemma}
	\label{lem:runningTime}
	There is a function $f:\mathbb{R} \rightarrow \mathbb{R}$ such that Algorithm~\ref{alg:EPTAS} runs in time $f(\frac{1}{\eps}) \cdot |\cm|^{O(1)}$. 
\end{lemma} The following is the main result of this paper. 

\begin{theorem}
	Algorithm~\ref{alg:EPTAS} is an EPTAS for BMI.
\end{theorem}

\begin{proof}
	Let $\cm =  (E,\cI,s,p)$ be a BMI instance and let $0<\eps<1$. By Lemma~\ref{lem:runningTime}, the running time of Algorithm~\ref{alg:EPTAS} is of an EPTAS algorithm. In addition, by Step~\ref{step:init} and Step~\ref{step:update} of Algorithm~\ref{alg:EPTAS} the returned set $A$ is a solution of $\cm$. Therefore, by the above it suffices to show a lower bound on the profit of $A$. 
	
	Let $U$ be a solution of $\cm$ with profit $p(U) = \OPT(\cm)$. In the following we partition a subset (with possibly empty sets in the partition) of the profitable elements in $U$ into $\eps^{-1}$ disjoint sets. Formally, for all $i \in \left[\eps^{-1}\right]$ let 
	
	\begin{equation}
		\label{J_i}
J_i = \big\{\ell \in H \cap U~\big|~ p(\ell) \in \big[\eps^{i+1} \cdot p\left(B^*\right), \eps^{i} \cdot p\left(B^*\right) \big) \big\}		
	\end{equation}
	
 be all elements in $H \cap U$ with profits in the interval $\big[\eps^{i+1} \cdot p\left(B^*\right), \eps^{i} \cdot p\left(B^*\right) \big) $. Since $J_1,\ldots, J_{\eps^{-1}}$ are $\eps^{-1}$ disjoint  subsets of elements within $U$, there is $i^* \in \left[\eps^{-1}\right]$ such that $p(J_{i^*}) \leq \eps \cdot \OPT(\cm)$ as otherwise it implies that $p(U) > \OPT(\cm)$ in contradiction. Now, define $J_0 = (H \cap U) \setminus \bigcup_{i \in [\eps^{-1}]} J_i$ as the largest profitable elements in $H \cap U$ not counted in the above subsets; in addition, let $L =  \bigcup_{k < i^*} J_k$ (including $J_0$) be the {\em most profitable} and $Q = U \setminus (L \cup J_{i^*})$ the {\em less profitable} elements. 
 
 Since $L \subseteq H \cap U$, by the hereditary property and Lemma~\ref{q} it holds that $L \in \cI^q$; hence, by Theorem~\ref{lem:main} there is $Z \subseteq S$, where $S = \textsf{Chosen}(\cm,T,q)$, such that $s(Z) \leq s(L)$, for all $ y \in T$ it holds that $|Z \cap E_y| = |E_y \cap L|$ and $Z \in \cI^q$.  Then, because $Z \subseteq S$, there is an iteration of Algorithm~\ref{alg:EPTAS} such that $F = Z$ in Step~\ref{step:for}. We use the following auxiliary claim whose proof is given in Section~\ref{sec:proofs}.
	
	\begin{claim}
		\label{claim:x}
		There is a point $\bar{x}^*$ in the $Z$-LP such that $p(Z)+\sum_{\ell \in E_Z} \bar{x}^* \cdot p(\ell) \geq (1-3 \eps) \OPT(\cm)$.
	\end{claim}


	
	
	Recall the optimal vertex solution to the $Z$-LP, that is $\bar{\lambda}^Z$, constructed in Step~\ref{step:vertex} of Algorithm~\ref{alg:EPTAS} in the iteration for $F = Z$. Now, 
	
	
		\begin{equation*}
		\label{eq:p(x)}
		\begin{aligned}
			p(A) \geq p(C(Z)) \geq  p(Z)+\sum_{\ell \in E_Z} \bar{\lambda}^Z \cdot p(\ell) -2\eps \cdot p(B^*) \geq (1-3 \eps) \OPT(\cm) - 2\eps \cdot p(B^*) \geq  (1-5 \eps) \OPT(\cm)
		\end{aligned}
	\end{equation*} The first inequality is because $Z \in \cI$ and $s(Z) \leq s(L) \leq s(U)$; it follows that $C(Z)$ is a solution of $\cm$ and $A$ is updated in Step~\ref{step:update} to be a solution of at least the same profit. The second inequality is because \eqref{LP} is in a face of dimension $1$ of $P_{M_Z}$; then, by Lemma~\ref{lem:integral} there are at most $2$ fractional entries in $\bar{\lambda}^Z$; in addition, the fractional entries correspond to elements in $E_Z$ with profits at most $\eps p(B^*)$ by \eqref{E_F} and \eqref{J_i}. The third inequality is by Claim~\ref{claim:x} and because $\bar{\lambda}^Z$ is an optimal solution to the $Z$-LP. The last inequality is because $B^*$ is a solution of $\cm$. By the above and using parameter $\delta = \frac{\eps}{5}$, we have an EPTAS for BMI. 
	%(\max\{|Z|-|Q|,0\}) \eps^{i^*+1} \cdot p(B^*)  

\end{proof}

\comment{arg1
We construct the {\em profit classes} as the set containing the possible profits of profitable elements after the rounding (i.e., by $\bar{p}$), that is $T = \{\bar{p}(\ell)~|~\ell \in H\}$ and assume that for all $\ell \in I$ it holds that $p(\ell) \leq \OPT(\cI)$ since otherwise $\ell$ does not belong to any solution of $\mathcal{I}$.


\begin{lemma}
	\label{T}
	$|T| \leq q$. 
\end{lemma}

\begin{proof}
	
	
	Let $r^* =  \max_{\{r \in \mathbb{N}~|~ r \cdot \eps \cdot h \in T\}} r$. If $r^* >q$, it holds that there is an element$\ell \in E$ such that $$p(\ell) = r^* \cdot \eps \cdot h > q \cdot \eps \cdot  h = 2 \cdot \eps^{-t-1} \cdot \eps \cdot \eps^{t} \cdot p(B^*) = 2p(B^*)  \geq \OPT(\cm).$$ We reach a contradiction since for all $\ell \in E$ it holds that $p(\ell) \leq \OPT(\cm)$ and the claim follows. 
	
\end{proof}
}


\comment{



In the following we describe the construction of a subset of elements $S \subseteq I$ for the proof of Theorem~\ref{thm:guess}. Let $y \in T$  be a profit class. We choose up to $2q$ groups, and from each group we take up to $2q$ elements with profit $y$ for $S$. This is done greedily, where we take elements small as possible as explained next more concretely.  

We define the following order on $\cG$. Let $y \in T$ be a profit class. Let $\cG_y = \{G \in \cG~|~\exists \ell \in G \text{ s.t. } p(\ell) = y\}$. For all $G \in \cG_y$, let \begin{equation}
	f_y(G) = \min_{\ell \in G \text{ s.t. } p(\ell) = y} s(\ell).
\end{equation} In words, $f_y(G)$ is the minimal size of an elementin $G$ with profit exactly $y$.  Let $|\cG_y| = n$ and let $G_1, \ldots, G_n$ be a non decreasing order of $\cG_y$ by the function $f_y$. Finally, let $q_y = \max\{2q, |\cG_y|\}$. These are the {\em chosen groups} for $y$ and let $F(y) = \{G_1, \ldots, G_{q_y}\}$ be this set of groups. For each chosen group $G \in F(y)$ we choose the {\em chosen elements} as follows. Let $G(y) = \{\ell \in G~|~p(\ell) = y\}$ be the elements in $G$ with profit $y$. Now, let $\ell_1, \ldots, \ell_{\alpha}$ be the elements in $G(y)$ according to a non decreasing order by the elementsizes (the function $s$). We define the chosen elements as the set $G^i(y) =  \{\ell_1, \ldots, \ell_{\beta}\}$ where $\beta = \min\{2q, |\cG(y)|\}$ elements in the above order. 

Finally, let the $y${\em-elements} be the union of the chosen elements for all chosen groups of $y$:  $S_y = \bigcup_{i \in [q_y]} G^i(y)$. Then, our algorithm then computes the union of the $y$-elements for all profit classes $y \in T$. That is, $S = \bigcup_{y \in T} S_y$. The pseudocode of Algorithm \textsf{Profitable} is given in Algorithm~\ref{alg:guess}. 


\comment{arg1
	\begin{equation}
		\label{D}
		S_y = \bigcup_{i \in [q_y]} G^i(y). 
\end{equation}} 


\begin{algorithm}[h]
	\caption{$\textsf{Profitable}(\mathcal{I} = (I,\cG,s,p,m,c), \eps)$}
	\label{alg:guess}
	%Compute the instance $\mathcal{J} = (I,\cG,s,\bar{p},m,c)$. 
	
	For each $y \in T$ compute the $y$-elements $S_y$. 
	
	Return $S = \bigcup_{y \in T} S_y$. 
\end{algorithm}


\begin{lemma}
	\label{S}
	$S$ returned by Algorithm~\ref{alg:guess} satisfies $|S| \leq \exp(\eps^{-3})$.
\end{lemma}

\begin{proof}
	We use the following inequality:
	
	\begin{equation}
		\label{q3}
		\begin{aligned}
			q^3 ={} & \left(\ceil{2 \cdot \eps^{-t-1}} \right)^3 \leq \left(2 \cdot \eps^{-t-1}+1\right)^3 \leq \left(2 \cdot \eps^{-t-1}+2 \cdot \eps^{-t-1}\right)^3 \leq  \left(\eps^{-t-2}\right)^3 \\ 
			\leq{} &	\left(\eps^{-\eps^{-1}-3}\right)^3 \leq 	\left(\eps^{-2\eps^{-1}-}\right)^3 = \eps^{-6\eps^{-1}} \leq \eps^{-\eps^{-2}}   \leq \exp(\eps^{-1})^{\eps^{-2}} \leq \exp(\eps^{-3}). 
		\end{aligned}	
	\end{equation} The inequalities follow since $0<\eps<0.1$. Now, 
	
	$$|S| = \left|\bigcup_{y \in T} S_y\right| \leq \bigcup_{y \in T} \left|S_y\right| \leq |T| \cdot 2q\cdot 2q \leq 4q^3 \leq 4 \exp(\eps^{-3}) \leq \exp(\eps^{-4}).$$ The second inequality is by the definition of the $y$-elements. The third inequality is by Lemma~\ref{T}. The fourth inequality is by \eqref{q3}. 
\end{proof}

\begin{lemma}
	The running time of Algorithm~\ref{alg:guess} is $f(1/\eps) \cdot |\cI|^{O(1)}$. 
\end{lemma}

\begin{proof}
	Computing $X$ is polynomial in $ |\cI|^{O(1)}$ by \cite{fleischer2011tight}. Then, we can also compute $\cJ$ in polynomial time. As $|T|$ and $q$ are bounded by a function of $1/\eps$, we can also compute the $y$-elements for all $y \in T$ in polynomial time. 
\end{proof}

For the proof of Theorem~\ref{thm:guess}, given an optimal solution $Z$ of $\mathcal{J}$ we construct a solution as explained next. This requires some definitions. Let $A$ be a solution of $\mathcal{J}$ and let the profitable elements assigned to a bin of $A$ be $$\textsf{elements}(A) = H \cap \bigcup_{i \in [m]} A_i.$$ Moreover, let $\cG_A = \{G \in \cG~|~ G \cap \textsf{elements}(A) = \emptyset\}$ be the set of groups that do not contain a profitable elementin the solution $A$. For all $\ell \in I$, $G \in \cG$ such that $\ell \in G$, let $\textsf{group}(\ell) = G$ be the {\em group of} $\ell$. Finally, for all $y \in T$ let $S_y[A] = \{\ell \in S_y ~|~ \textsf{group}(\ell) \notin \cG_A\}$ be the elements in $S_y$ of groups disjoint to the solution $A$ and let the elementof smallest size in this set be\footnote{assume that the elements have distinct sizes, or choose arbitrarily an elementof smallest size.}

\begin{equation}
	\label{item}
	\ell_y[A] = \arg \min_{\ell \in S_y[A]} s(\ell). 
\end{equation} Furthermore, let $A_y = \textsf{elements}(A) \cap \{\ell \in I ~|~p(\ell) = y\}$ be profitable elements in $A$ with profit $y$. We do the following process. We start with an optimal solution $Z^0$ for $\mathcal{J}$. If there is $y \in T$ such that $Z^0_y \setminus S \neq \emptyset$, then let  $i \in [m]$ such that there is $\ell \in (Z^0_i \cap Z^0_y) \setminus S$; we define an {\em exchange} in $Z_i$: let $Z^1_i = (Z^0_i \setminus \{\ell\}) \cup \{\ell_y[Z^0]\}$ and for all $j \in [m] \setminus \{i\}$ let $Z^1_j = Z^0_j$; finally, let $Z^1 = (Z^1_1, \ldots, Z^1_m)$ be the solution after one iteration and let $S^1 = S \setminus \{\ell_y[Z^0]\}$.

In general, after $r$ iterations, the process finds $y \in T$ (if exists) such that $Z^r_y \setminus S^{r} \neq \emptyset$, then chooses arbitrary $i \in [m]$ such that there is $\ell \in (Z^{r}_i \cap Z^{r}_y) \setminus S^{r}$; we define an {\em exchange} in $Z^r_i$: let $Z^{r+1}_i = (Z^{r}_i \setminus \{\ell\}) \cup \{\ell_y[Z^{r}]\}$ and for all $j \in [m] \setminus \{i\}$ let $Z^{r+1}_j = Z^r_j$; finally, let $Z^{r+1} = (Z^{r+1}_1, \ldots, Z^{r+1}_m)$ be the solution after iteration $r$. The pseudocode of the process is given in Algorithm~\ref{alg:2}. 

\begin{algorithm}[h]
	\caption{$\textsf{solution}(\cJ, Z^0)$}
	\label{alg:2}
	
	Initialize $r = 0$. 
	
	
	\While{$\exists y \in T$ such that $Z^r_y \setminus S^{r} \neq \emptyset$ }{
		
		Let $i \in [m]$ such that there is $\ell \in (Z^{r}_i \cap Z^{r}_y) \setminus S^{r}$. 
		
		Let $Z^{r+1}_i = (Z^{r}_i \setminus \{\ell\}) \cup \{\ell_y[Z^{r}]\}$ and for all $j \in [m] \setminus \{i\}$ let $Z^{r+1}_j = Z^r_j$. 
		
		Let $Z^{r+1} = (Z^{r+1}_1, \ldots, Z^{r+1}_m)$
		
		Update $r \leftarrow r+1$. 
	}
	
	Return $B = Z^r$ and $A = (Z^r_1 \cap H, \ldots, Z^r_b \cap H)$. 
	
\end{algorithm}

%While there are $i \in [m]$, $\ell \in Z_i \setminus S$ such that $p(\ell) = y$, then let 

%We use the following bin structuring technique from \cite{jansen2010parameterized}. 
}

\section{Properties of the chosen elements}
\label{sec:lemMainProof}

In this section we prove Theorem~\ref{lem:main} mainly based on simple properties of matroids. %For all $y \in T$, By the following lemma, given some independent set $F$, each elementin $F$ can be replaced by an elementfrom the chosen elements, with the same profit and smaller size, while resulting in an independent set.  

\begin{lemma}
	\label{lem:1}
	For all $y \in T$, $F \in \cI^q$, and $d \in F \cap \left(E_y \setminus B_y \right)$, there is $b \in B_y$ such that (i) $F-d+b \in \cI^q$ and (ii) s$\left(b\right) \leq s(d)$. 
\end{lemma}


\comment{ \begin{equation}
		\label{R}
		R(i) = \{b_j~|~i \in [n], b_j < b_i\}
\end{equation}}

\begin{proof}
\comment{arg1	Let $n = |B_y|$ and for all $i \in [n]$ let $b_i$ be the $i$-th elementadded to $B_y$ in Algorithm \textsf{Chosen}; in addition, let $	R(i) = \{b_j~|~i \in [n], b_j < b_i\}$ be the first elements added to $B_y$ until the predecessor elementto $y(b_i)$.

 Let the elementadded to $B_y$ which satisfies (i) and with minimal size be \begin{equation}
		\label{b*}
		b^* = \arg \min_{\{b \in B_y ~|~  F-d+b \in \cI\}} s(b).
	\end{equation} If there is more than one such $b^*$ choose the first such elementadded to $B_y$. We use the following variable. First, if there is no $b \in B_y$ such that $F-d+b \in \cI$ then let $R = \emptyset$; otherwise, let $R = \{b \in B_y~|~ s(b)<s(b^*)\}$ be all elements added to $B_y$ before element$b^*$. Observe that for all $b \in R$ it holds that $s(b)<s(b^*)$ since $b^*$ is the first elementadded to $B_y$ with size $s(b^*)$. Assume towards a contradiction that $s\left(b^*\right) > s(d)$. We reach in two complementary cases, as follows. }


Define the subset of elements in $B_y$ satisfying (i) and (ii) as: 

\begin{equation}
	\label{R}
	R = \{b \in B_y~|~ F-d+b \in \cI^q, s(b) \leq s(d)\}. 
\end{equation} If $R \neq \emptyset$ the claim follows by \eqref{R}. Therefore, Assume towards a contradiction that $R = \emptyset$. Because $d \in F$, it follows that $|R|<|F|$. Then, by Property~1 of Lemma~\ref{lem:prel} there is $C \in F \setminus R, |C| = |F|-|R|$ such that $R \cup C \in \cI^q$. We use the following auxiliary claim. 

\begin{claim}
\label{dC}
$d \notin C$. 
\end{claim}
%Second, $|B| \leq |R|<|F| \leq q$ (recall that $F \in \cI^q$) and it holds that $|B| \leq q-1$. 
\begin{proof}
Assume towards a contradiction that $d \in C$; thus, since $R \cup C \in \cI^q$ and $\{d\} \subseteq C$, we get by the hereditary property that $R+d \in \cI^q$. Let $B$ be the temporary set $B_y$ right before Step~\ref{step:if} of Algorithm~\ref{alg:chosen} for $y(i) = d$. Observe that $B \subseteq R$ by \eqref{R} and by the ordering of the elements in Step~\ref{order} of Algorithm~\ref{alg:chosen}. Then, because $R+d \in \cI^q$ then by the hereditary property it follows that $B+d \in \cI^q$. Therefore, by Step~\ref{step:if} of Algorithm~\ref{alg:chosen}, we should add $d$ to $B_y$ and it follows that $d \in B_y$ in contradiction. 
\end{proof} 	Therefore, $|R \cup C| = |F|> |F-d|$ since $d \in F$; by the exchange property of $(E,\cI^q)$, there is $b^* \in (R \cup C) \setminus (F-d)$ such that $F-d+b \in \cI^q$. Since $C \subseteq F$ and $d \notin C$ by Claim~\ref{dC}, it holds that $C \subseteq F-d$; it follows that $b^* \in R$ in contradiction that $R = \emptyset$. 
%Observe that $d \notin R$ since $d \notin B_y$. 
\end{proof}

%We split the proof into two complementary cases and in each show the existence of $b^* \in B_y$ which satisfies (i) and (ii). 
	\comment{arg1
	\begin{itemize}
		\itemIf $|R| \geq |F|$. Then, it follows that $|R| \geq |F| > |F-d|$ since $d \in F$. Therefore, by the exchange property of $(E,\cI^q)$, there is $b^* \in R \setminus (F-d)$ such that $ F-d+b^* \in \cI^q$. Hence, by \eqref{R} we have that $b^*$ satisfies the conditions of the lemma.%By the definition of $R$ we have that $s(b)<s(b^*)$ in contradiction to \eqref{b*}.
		
		\itemIf $|R| < |F|$. 
	}
		
		


	\noindent{\bf Proof of Lemma~\ref{lem:main}}:  Let $Z \in \cI^q$ such that (a) $s(Z) \leq s(G)$, (b)  for all $y \in T$ it holds that $|E_y \cap Z| = |E_y \cap G|$ and $|Z \cap S|$ is maximal among all $Z' \in \cI^q$ satisfying (a) and (b). Assume towards a contradiction that there is $a \in Z \setminus S$ and let $y = p(a)$.  Because $a \notin S$ and $B_y \subseteq S$, by Lemma~\ref{lem:1} there is $b^* \in B_y$ such that (i) $Z-a+b^* \in \cI^q$ and (ii) $s(b^*) \leq s(a)$; we define $W = Z-a+b^*$. It holds that $s(W) = s(Z)-s(a)+s(b^*) \leq s(Z)$ by (ii). In addition, for all $y' \in T -y$ it holds that $|E_{y'} \cap W| = |E_{y'} \cap G|$ and $|E_{y} \cap W| = |E_{y} \cap G|$ because $a, b^* \in E_y$ and by (b). Furthermore, $|W \cap S| > |Z \cap S|$ because $a \notin S$ and $b^* \in S$. Hence, $W$ satisfies (a), (b) and $|W \cap S| > |Z \cap S|$ which contradicts the maximality of $|Z \cap S|$. \qed

\section{Deferred Proofs}
\label{sec:proofs}


	\noindent{\bf Proof of Lemma~\ref{lem:J}}: 
	We use the following inequality. For any $\ell \in H$, by \eqref{eq:max} there is $r \in \mathbb{N}$ such that $\bar{p}(\ell) = r \cdot \eps \cdot h$ and $p(\ell) < (r+1) \cdot \eps \cdot h$; therefore,  		\begin{equation}
			\label{claim:1}
			p(\ell)-\bar{p}(\ell) \leq (r+1) \cdot \eps \cdot h- r \cdot \eps \cdot h =  \eps \cdot h \leq \eps \cdot p(\ell).
		\end{equation} The first inequality is by \eqref{eq:max}. The second inequality is because $\ell \in H$. Now, 
%	\end{proof} 


	
	$$\bar{p}(A) = \sum_{\ell \in A} \bar{p}(\ell) \geq \sum_{\ell \in A} (1-\eps) \cdot p(\ell) = (1-\eps) \sum_{\ell \in A} p(\ell) = (1-\eps) \cdot p(A).$$ The inequality follows for all $\ell \in E \setminus H$ since $\bar{p}(\ell) = p(\ell) \geq (1-\eps) p(\ell)$; furthermore, for all $\ell \in H$ the inequality follows by \eqref{claim:1}.\qed






	\noindent{\bf Proof of Lemma~\ref{q}}: 
	We split the proof into the two parts (i) and (ii). 

\begin{itemize}
	\item(i). 	Assume towards a contradiction that there is a solution $A$ of $\cm$ such that $q < |H \cap A|$. Therefore, 
	
	
	
	\begin{equation}
		\label{contradiction1}
		\begin{aligned}
			p(A) ={} & \sum_{\ell \in A} p(\ell) \geq  \sum_{\ell \in H \cap A} p(\ell) \geq \sum_{\ell \in H \cap A} h   = h \cdot |H \cap A| > h \cdot q = \eps^{t} \cdot p\left(B^*\right) \cdot 2 \eps^{-t-1} \\ ={} & 2 \cdot p\left(B^*\right)  \eps^{-1} \geq \OPT(\cm) \cdot \eps^{-1}  >   \OPT(\cm). 
		\end{aligned}
	\end{equation} The second inequality is because for all $\ell \in H$ it holds that $p(\ell) \geq h$. The third inequality is by the assumption that $q < |H \cap A|$. the fourth inequality is because $B^*$ is a $2$-approximation for $\cm$. By \eqref{contradiction1} we reach a contradiction that $A$ is a solution of $\cm$ with profit strictly larger than $\OPT(\cm)$. 
	
	\item(ii). 	Let $r^* =  \max_{\{r \in \mathbb{N}~|~ r \cdot \eps \cdot h \in T\}} r$. If $r^* >q$, it holds that there is an element$\ell \in E$ such that $$p(\ell) = r^* \cdot \eps \cdot h > q \cdot \eps \cdot  h = 2 \cdot \eps^{-t-1} \cdot \eps \cdot \eps^{t} \cdot p(B^*) = 2p(B^*)  \geq \OPT(\cm).$$ We reach a contradiction since $A = \{\ell\}$ is a solution with $p(A) > \OPT(\cm)$. By the above, $|T| \leq |\{r \in[q]\}| = q.$ \qed
\end{itemize}

	\noindent{\bf Proof of Lemma~\ref{lem:runningTime}}: 
	
		\noindent{\bf Proof of Claim~\ref{claim:x}}: 
		Because $U \in \cI$ and $Q \subseteq U$, by the hereditary property it follows that $Q \in \cI$. We define a point in the $Z$-LP relied on the following defined subset of elements. First, if $|Z| > |Q|$ then define $C = \emptyset$. Otherwise, $|Q| \geq |Z|$; then, by Lemma~\ref{lem:prel} there is $C' \subseteq Q \setminus Z, |C'|  = |Q| - |Z|$ such that $Z \cup C' \in \cI$; in this case, define $C = C'$. Now, define $\bar{x}^* \in \{0,1\}^{E_Z}$ such that for all $\ell \in C$ define $\bar{x}^*_{\ell} = 1$ and for all $\ell \in E_Z \setminus C$ define $\bar{x}^*_{\ell} = 0$. It is shown below that $\bar{x}^*$ is in the $Z$-LP. Because $Z \cup C \in \cI$, it follows by the definition of the matroid polytope that $\bar{x}^* \in P_{M_Z}$, the matroid polytope of $Z$. In addition, $$\sum_{\ell \in E_Z} \bar{x}^* \cdot s(\ell) = \sum_{\ell \in C} s(\ell) \leq \sum_{\ell \in Q} s(\ell) = s(Q) \leq 1-s(U \setminus Q) \leq 1-s(L) \leq 1-s(Z).$$ By the above and by \eqref{LP}, it holds that $\bar{x}^*$ is in the $Z$-LP. For the profit, we use the following inequalities.
		
		\begin{equation}
			\label{eq:11}
			\begin{aligned}
				\eps \cdot p(Z)-p(Q \setminus C) \geq \eps \cdot p(Z)- |Z| \cdot \eps \cdot  \min_{\ell \in Z} p(\ell) \geq \eps \cdot p(Z)-\eps \cdot p(Z) = 0.
			\end{aligned}
		\end{equation} The first inequality is because $C \subseteq Q \setminus Z$, $|C|  = |Q| - |Z|$ or $C = \emptyset$ if $|Z| > |Q|$ and it follows that $|Q \setminus C| \leq |Z|$; in addition, by the definition of $E_Z$ all elements in $Q \setminus C$ are of profits at most $\eps \cdot \min_{\ell \in Z} p(\ell)$. In addition, 
		
		\begin{equation}
			\label{eq:p(x)2}
			\begin{aligned}
				p(Z)+p(Q) \geq{} & \bar{p}(Z)+p(Q)  =  (\bar{p}(L)+p(Q)) \geq  (1-\eps)\cdot p(L)+p(Q) \geq  (1-\eps) \left( p(L)+p(Q)\right) \\ = {} &  (1-\eps)  p(U \setminus J_{i^*}) \geq (1-\eps) (1-\eps)\OPT(\cm) \geq (1-2\eps)  \OPT(\cm). 
			\end{aligned}
		\end{equation} The first equality is because $Z,L \subseteq H$ and for all $y \in T$ it holds that $|Z \cap E_y| = |E_y \cap L|$. The second inequality is by Lemma~\ref{lem:J}. Therefore,
		
		\begin{equation*}
			\label{eq:p(x)}
			\begin{aligned}
				p(Z)+\sum_{\ell \in E_Z} \bar{x}^* \cdot p(\ell) ={} & p(Z)+\sum_{\ell \in C} p(\ell) = p(Z)+p(C) = p(Z)+p(Q) - p(Q \setminus C) \\ \geq{} & (1-\eps) (p(Z)+p(Q))+\eps \cdot p(Z)-p(Q \setminus C)  \geq (1-3\eps) \OPT(\cm).
			\end{aligned}
		\end{equation*} The last inequality is by \eqref{eq:11} and \eqref{eq:p(x)2}.  Hence, $\bar{x}^*$ satisfies the conditions of the claim. \qed